\documentclass[12pt,twoside,a4paper]{article} %12 punkters typsnitt (standard är 10), A4-format, och dokumenttyp är "article" jämför: report - längre rapporter med  kapitel, book - förberett för omslag osv., slides - så du kan göra presentationer, letter - för att skriva brev, IEEEtran - för journalartiklar (pröva för att se vad som händer, byt article mot IEEEtran)

\usepackage{geometry} % Möjliggör t.ex. A4

\usepackage{graphicx} % Required for inserting images

\usepackage{hyperref} %möjliggör länkar

\usepackage[utf8]{inputenc} %möjliggör olika teckenkodningssystem
\usepackage[swedish]{babel}

\usepackage{lineno} % möjliggör randnumrering för granskning!
\linenumbers % stänger av och på radnumrering

\title{DA211A}
\author{S. T. Udent, TGxxx, DA211A}
\date{September 2024} % Om du vill ha automatiskt datumsättning så kommentera bort den här raden

\begin{document}

%-tecknet betyder att man kommenterar bort raden

\maketitle % \maketitle sätter in fält såsom t.ex. \author och \date och \title naturligtvis
%-----------------------Här börjar min text ----------------------------




\section{Introduktion}
En algoritm är en sekvens av steg som garanterar en korrekt beräkning~\cite{kernighan2021}. % kursboken s. 80

\section{Ett nytt stycke}
Lite text
\subsection{En underrubrik}
och lite text här också

% referenslistan
% skriver vi långt så behövs en sidbrytning innan
%
ewpage % kommentera fram vid längre uppgifter
\bibliographystyle{ieeetr}
\bibliography{bibliography}
\end{document}% dokumentet avslutas